% !TeX root = ../偏微分方程笔记.tex
% 由主文档通过 \input 载入；请编译主文档。
\setcounter{chapter}{0} % 保留原书章号。
\chapter{微积分与常微分方程工具}
\label{app:calculus}

\section{多元微分与坐标变换}
若 $y=\Phi(x)$ 为 $C^1$ 映射，$u(x)=v(\Phi(x))$，则链式法则为 $\nabla_xu=D\Phi(x)^{\mathsf T}\nabla_yv$。若 $\det D\Phi(x_0)\ne0$，逆函数定理保证 $x_0$ 附近存在 $C^1$ 局部逆。非零雅可比给出的是局部结论；全局单射需要另证。

平面极坐标 $x=r\cos\theta$、$y=r\sin\theta$ 的面积元为 $r\,dr\,d\theta$，三维球坐标 $x=(r\sin\theta\cos\varphi,r\sin\theta\sin\varphi,r\cos\theta)$ 的体积元为 $r^2\sin\theta\,dr\,d\theta\,d\varphi$。一般 $n$ 维径向积分为
\[
 \int_{\mathbb R^n}f(|x|)\,dx=\sigma_{n-1}\int_0^\infty f(r)r^{n-1}\,dr,
\]
前提是被积函数非负或绝对可积。对径向函数 $u(x)=U(|x|)$，在 $r>0$ 有 $\Delta u=U''+(n-1)U'/r$；原点正则性必须由 $U$ 的极限另行检查。

\section{积分恒等式与边界方向}
若 $\Omega$ 有界且边界分片 $C^1$，向量场 $F\in C^1(\overline\Omega)$，则散度定理为
\[
 \int_\Omega\operatorname{div}F\,dx=\int_{\partial\Omega}F\cdot\nu\,dS.
\]
取 $F=u\nabla v$ 得式\eqref{eq:green-first}，交换变量再相减得式\eqref{eq:green-second}。若区域被挖去小球，内边界的区域外法向指向小球中心，这正是基本解计算中容易出错的符号。

区间 $[a,b]$ 的边界法向在左端为 $-1$、右端为 $1$，所以边界项 $[uv']_a^b=u(b)v'(b)-u(a)v'(a)$ 与多维记号完全一致。对较弱的函数空间，这些恒等式需通过稠密逼近和迹延拓后使用。

\section{常微分方程与参数依赖}
若 $F(t,z)$ 在开域内连续，并对 $z$ 局部利普希茨且常数在紧时间区间上一致，则初值问题 $z'=F(t,z)$、$z(t_0)=z_0$ 在该点附近存在唯一 $C^1$ 解。只要在有限时间内轨道留在定义域的某个紧子集中，就能继续延拓。若 $F$ 连续但不利普希茨，唯一性可能失效，例如 $z'=2\sqrt{|z|}$、$z(0)=0$ 可以有延迟离开零点的解。

线性方程 $y'+a(t)y=b(t)$ 的积分因子为 $e^{A(t)}$，$A(t)=\int_{t_0}^t a(s)ds$，因此
\[
 y(t)=e^{-A(t)}\left(y(t_0)+\int_{t_0}^t e^{A(s)}b(s)ds\right).
\]
若 $a,b$ 仅局部可积，此公式仍定义绝对连续解并使方程几乎处处成立。常系数振子 $y''+\omega^2y=b(t)$（$\omega>0$）则有
\[
 y(t)=y(0)\cos\omega t+\frac{y'(0)}\omega\sin\omega t
       +\int_0^t\frac{\sin\omega(t-s)}\omega b(s)\,ds.
\]
当 $\omega=0$，应取极限核 $t-s$。边值问题不同于初值问题，即使系数光滑，指定两端数据也可能无解或不唯一；谱问题正体现了这种区别。

